% Generated draft: requires mathematical and source review.
\documentclass[11pt]{article}
\usepackage[margin=25mm]{geometry}
\usepackage{fontspec}
\setmainfont{Latin Modern Roman}
\usepackage{amsmath,amssymb,mathtools,unicode-math}
\setmathfont{Latin Modern Math}
\usepackage{enumitem,needspace,xurl,hyperref}
\setlength{\parindent}{0pt}
\setlength{\parskip}{5pt}
\setlength{\emergencystretch}{3em}
\newcommand{\N}{\mathbb{N}}
\newcommand{\Z}{\mathbb{Z}}
\newcommand{\Q}{\mathbb{Q}}
\newcommand{\R}{\mathbb{R}}
\newcommand{\C}{\mathbb{C}}
\newcommand{\F}{\mathbb{F}}
\newcommand{\A}{\mathbb{A}}
\newcommand{\PP}{\mathbb{P}}
\newcommand{\E}{\mathbb{E}}
\newcommand{\eps}{\varepsilon}
\newcommand{\dd}{\,\mathrm{d}}
\begin{document}
\section*{UnsolvedMath ID: 30000844}
\textbf{Source problem code:} catalogue entry
\subsection{Definitions and mathematical statement}
The source asks whether the secondary $\rho$-invariants for Dirac operators in the noncommutative or foliated settings discussed in its report can yield geometric applications. A Dirac-type operator is a first-order differential operator whose principal symbol is Clifford multiplication; in a foliated setting a longitudinal such operator acts along the leaves. These terms identify the setting but not a particular secondary invariant. A $\rho$-invariant may refer to an eta-invariant difference, a measured invariant, or a higher class in a specified operator-algebraic target; the extract gives neither its construction nor its target and these alternatives are not equated. The question retained here is therefore the report's research direction: can its particular secondary $\rho$-invariants be used to deduce geometric consequences? No definite class of geometric objects, equivalence relation, obstruction, rigidity assertion, or nonvanishing condition is supplied, so there is no uniquely determined quantified theorem to substitute for this question. The clean/current version says new geometric applications, whereas the original asks for geometric applications in the settings discussed in the report; neither the intended application nor the reference point for novelty is stated. The supplied 2026-08-21 triage labels the record open, calls the invariant unspecified in the extraction, and reports the passage as work in progress with no uniquely testable resolution identifiable.
\subsection{Short English statement}
The report asks for geometric applications of its secondary Dirac $\rho$-invariants. The extract does not identify the invariant, its target, or the desired geometric consequence, so this remains a source-specific research direction rather than a complete conjecture.
\subsection{Sources}
\begin{itemize}\item Local UnsolvedMath registry, permanent ID 30000844, source problem catalogue entry, \emph{Geometric Applications of Secondary Rho Invariants}, proposed year 2007; no individual proposer supplied.\item Supplied citation: \emph{Noncommutative Geometry} (2008), Oberwolfach Reports, DOI 10.4171/owr/2007/43, pp.\ 2543--2608: \url{https://doi.org/10.4171/owr/2007/43}. Supplied literature label: original Oberwolfach report.\end{itemize}
\end{document}
